\documentclass[11pt]{jlreq}
\usepackage{amsmath, amssymb, amsfonts}
\usepackage{geometry}
\usepackage{fancyhdr}

% 余白設定
\geometry{left=20mm, right=20mm, top=25mm, bottom=25mm}

% 【追加修正】ヘッダーの高さ不足警告を解消する設定
\setlength{\headheight}{20pt}

% 2ページ目以降のデフォルト設定を「plain（ヘッダーなし、ページ番号のみ）」にする
\pagestyle{plain}

% ヘッダー設定
\fancyhf{}
\cfoot{\thepage}

\begin{document}

\thispagestyle{fancy}
\lhead{2025年度 数学 学年末試験 解答用紙No.1}

\rhead{\underline{\hspace{1cm}}組\underline{\hspace{1cm}}番　氏名：\underline{\hspace{5cm}}}

\section*{[1] 必答問題（解答のみ記入すること）}

\begin{tabular}{|p{0.45\textwidth}|p{0.45\textwidth}|}
\hline
(1) & (2) \\ [5ex] \hline
(3) & (4) \\ [5ex] \hline
(5) & (6) \\ [5ex] \hline
(7) & (8) \\ [5ex] \hline
(9) & (10) \\ [5ex] \hline
\end{tabular}

\vspace{2em}

\hrule

\vspace{2em}

\section*{[2] 選択問題}

\noindent

\textbf{選択した問題に○をつける： [ \quad 2a \quad $\cdot$ \quad 2b \quad $\cdot$ \quad 2c \quad ]} \\

\subsection*{(1)}

\vfill

\newpage

\subsection*{(2)}

\vfill

\newpage

\thispagestyle{fancy}

\lhead{2025年度 数学 学年末試験 解答用紙No.2}

\rhead{\underline{\hspace{1cm}}組\underline{\hspace{1cm}}番　氏名：\underline{\hspace{5cm}}}

\subsection*{(3)}

\vfill

\newpage

\subsection*{(4)}

\vfill

\newpage

\end{document}